\section{Proofs of the operational and new results}\label{supp:proofs70}
\subsection{Proof of the result \ref{thm:screen67}}\label{supp:proof70-operational67-1}
\begin{proof}
Strong convexity gives $f((1-s)l+su)\le(1-s)f(l)+sf(u)-\mu s(1-s)(u-l)^2/2$. Replacing the endpoint values by their upper bounds proves \eqref{eq:screen67}. Subtraction of a valid lower bound proves \eqref{eq:screen-stop67}. A failed screen adds no endpoint to the partition and alters no lower bound. Subsequent protected refinement is therefore covered by the earlier termination proof. A successful screen supplies exactly the local accuracy premise used in the acquired-policy recursion. Finally, $C$ is a convex quadratic in $a$, proving the sublevel-set assertion.
\end{proof}

\subsection{Proof of the result \ref{thm:work-loss67}}\label{supp:proof70-operational67-2}
\begin{proof}
A successful screen returns without an endpoint query. A failed screen invokes a procedure bounded by $K(z)$, and its margin loss equals one. This proves \eqref{eq:work-bound67}. For each fixed predictor the variable $K(z)\ell_\tau(z,\theta)$ lies in $[0,K_{\max}]$. The bounded-sum inequality of \citet{hoeffding1963} and a union bound over the conditionally fixed family give \eqref{eq:work-validation67}.
\end{proof}

\subsection{Proof of the result \ref{prop:readout67}}\label{supp:proof70-operational67-3}
\begin{proof}
The chord is a convex quadratic in its normalized action, which is affine in $w$. Squaring the positive part preserves convexity and gives a continuously differentiable function. The nonnegative weighted sum and quadratic penalty preserve these properties. Convexity gives $F(w)-F(v)\le g^\top(w-v)$, and the minimum of $g^\top v$ on the box is $\sum_j\min\{0,g_j\}$. This proves \eqref{eq:dual-gap67}. Every failed screen has squared positive-part residual at least $\tau^2$, while every successful screen uses zero further endpoint queries. Summing the corresponding $K_i$ bounds proves the work claim. Strong convexity with modulus $\lambda$ proves the distance bound.
\end{proof}

\subsection{Proof of the result \ref{thm:vector67}}\label{supp:proof70-vector67-1}
\begin{proof}
Semiconcavity means that $f(a)-H\|a\|_2^2/2$ is concave. Its barycentric inequality, with vertex values replaced by lower endpoints, gives \eqref{eq:triangle67}. The tangent inequality for the convex quadratic $p$ proves \eqref{eq:dual-triangle67}; a linear function attains its simplex minimum at a vertex. Taking minima over a cover preserves the lower bound. Each feasible upper witness is an action of the actual model. The local gap premise therefore implies the earlier policy recursion, irrespective of how the triangle or support point was proposed.

For finiteness, use simplex lattice denominator $A$, with $(A+1)(A+2)/2$ vertices. On triangles of diameter proportional to $c(x)/A$, the curvature error tends to zero as $A^{-2}$. The endpoint error is reduced by finite descendant accuracy and quadrature refinements. The retained exact-rational recovery contract requires $\delta>16\Lambda\,2^{-32}$ at each recursive call, leaving positive slack after rounding; dyadic enlargement of $A$ and $q$ then terminates. This is a finite arithmetic procedure, not a uniform bound on rational bit cost. For \eqref{eq:two-shock67}, the independent innovation variances after conditioning on the bins are $\rho_j^2/(3q^2)$, their direction norms squared are $d$, and the semiconcavity Taylor remainder is $M_{t+1}/2$ times their sum. Substitution of $\rho_1=1/32$, $\rho_2=1/64$ gives the coefficient. Finally, each action query has $q_r^2$ descendants, proving \eqref{eq:vector-work67}.
\end{proof}

\subsection{Proof of the result \ref{thm:centered69}}\label{supp:proof70-centered69-1}
\begin{proof}
Conditional expectation and $b_{t+1}=b_t\beta_t$ imply
\[
 \mathbb E\sum_{t<T}b_t a_t^\pi(X_t^\pi)
 =\mathbb E\left[\sum_{t<T}b_tc_t(X_t^\pi,\pi_t(X_t^\pi))
                  +b_Tg(X_T^\pi)-V_0(X_0^\pi)\right].
\]
The intermediate value terms cancel. Subtracting the two expectations cancels the initial value under the common initial law, proving the identity. Summing the pointwise local bounds proves the range. Applying the bounded-sum inequality to each lower and upper endpoint sample, with failure allocation $\alpha/(2M)$, and taking a union bound proves \eqref{eq:centered-ci69}. Endpoint containment supplies the comparison between endpoint expectations and the target. The width identity follows by subtraction. No pathwise identity between $D$ and the realized cumulative-cost difference is asserted.
\end{proof}

\subsection{Proof of the result \ref{cor:precision69}}\label{supp:proof70-centered69-2}
\begin{proof}
The sample-size condition implies $h_n\le s$. Insert this and the numerical-width condition into the width identity. The final claims are set inclusions for the simultaneous confidence interval.
\end{proof}

\subsection{Proof of the result \ref{thm:deployment70}}\label{supp:proof70-deployment70-1}
\begin{proof}
The bounded-sum argument in Theorem~\ref{thm:work-loss67} gives $P f_\theta\le u_\theta$ for every $\theta$ on one event of probability at least $1-\alpha$. For $0\le f\le K_{\max}$, the layer-cake representation gives $|Pf-Qf|\le K_{\max}\|P-Q\|_{\rm TV}$. Hence, on that event,
\[
 \mathbb E[N_{\theta_b}(Z_b)\mid\mathcal F_{b-1}]
 \le\min\{K_{\max},u_{\theta_b}+K_{\max}\rho_b\}.
\]
The variables $c_bN_{\theta_b}(Z_b)$ have conditional ranges of length at most $c_bK_{\max}$. Conditional Hoeffding bounds, multiplied successively along the filtration and followed by the exponential Markov inequality, show that the sum exceeds its conditional means by at most the final term of \eqref{eq:deployment-work70}, except on an event of probability $\gamma$. Union with the validation failure event proves the result. Every selected readout remains a proposal to the same verified gate. Its statewise certificate, and thus the earlier acquired-policy recursion, is unchanged.
\end{proof}

\subsection{Proof of the result \ref{prop:adoption70}}\label{supp:proof70-deployment70-2}
\begin{proof}
The expected cost difference in a normal block is $m\Delta$; it is zero in a declined block. Linearity of expectation gives $m(B-D)\Delta$, without requiring block independence. Multiply the cost interval by its nonnegative price, and add the work difference and fee. All price choices use the same event, so no additional probability allocation is needed.
\end{proof}

